\section{USAMO 1982}
        \begin{enumerate}
        	\item (\textit{USAMO 1982}) In a party with $1982$ persons, among any group of four there is at least one person who knows each of the other three. What is the minimum number of people in the party who know everyone else?


 



\textbf{Solution}

We induct on $n$ to prove that in a party with $n$ people, there must be at least $(n-3)$ people who know everyone else. (Clearly this is achievable by having everyone know everyone else except three people $A, B, C$, who do not know each other.)

 Base case: $n = 4$ is obvious.

 Inductive step: Suppose in a party with $k$ people (with $k \ge 4$), at least $(k-3)$ people know everyone else. Consider a party with $(k+1)$ people. Take $k$ of the people (leaving another person, $A$, out) and apply the inductive step to conclude that at least $(k-3)$ people know everyone else in the $k$-person group, $G$.

 Now suppose that everyone in the group $G$ knows each other. Then take $3$ of these people and $A$ to deduce that $A$ knows a person $B \in G$, which means $B$ knows everyone else. Then apply the inductive step on the remaining $k$ people (excluding $B$) to find $(k-3)$ people out of them that know everyone else (including $B$, of course). Then these $(k-3)$ people and $B$, which enumerate $(k-2)$ people, know everyone else.

 Suppose that there exist two people $B, C \in G$ who do not know each other. Because $k-3 \ge 1$, there exist at least one person in $G$, person $D$, who knows everyone else in $G$. Now, take $A, B, C, D$ and observe that because $B, C$ do not know each other, either $A$ or $D$ knows everyone else of $A, B, C, D$ (by the problem condition), so in particular $A$ and $D$ know each other. Then apply the inductive step on the remaining $k$ people (excluding $D$) to find $(k-3)$ people out of them that know everyone else (including $D$, of course). Then these $(k-3)$ people and $D$, which enumerate $(k-2)$ people, know everyone else.

 This completes the inductive step and thus the proof of this stronger result, which easily implies that at least $1982 - 3 = \boxed{1979}$ people know everyone else.

 


	\item (\textit{USAMO 1982}) Let $S_r=x^r+y^r+z^r$ with $x,y,z$ real. It is known that if $S_1=0$,


 $(*)$ $\frac{S_{m+n}}{m+n}=\frac{S_m}{m}\frac{S_n}{n}$


 for $(m,n)=(2,3),(3,2),(2,5)$, or $(5,2)$. Determine \textit{all} other pairs of integers $(m,n)$ if any, so that $(*)$ holds for all real numbers $x,y,z$ such that $S_1=0$.


 



\textbf{Solution 1}

Claim Both $m,n$ can not be even.

 Proof$x+y+z=0$ ,$\implies x=-(y+z)$.

 Since $\frac{S_{m+n}}{m+n} = \frac{S_m S_n}{mn}$,

 by equating cofficient of $y^{m+n}$ on LHS and RHS ,get

 $\frac{2}{m+n}=\frac{4}{mn}$.

 $\implies  \frac{m}{2} + \frac {n}{2} = \frac{m\cdot n}{2\cdot2}$.

 So we have, $\frac{m}{2} \biggm{|} \frac{n}{2}$ and $\frac{n}{2} \biggm{|} \frac{m}{2}$.

 $\implies m=n=4$.

 So we have $S_8=2(S_4)^2$.

 Now since it will true for all real $x,y,z,x+y+z=0$.So choose $x=1,y=-1,z=0$.

 $S_8=2$ and $S_4=2$ so $S_8 \neq 2 S_4^2$.

 This is contradiction. So, at least one of $m,n$ must be odd. WLOG assume $n$ is odd and m is even. The coefficient of $y^{m+n-1}$ in $\frac{S_{m+n}}{m+n}$ is $\frac{\binom{m+n}{1} }{m+n} =1$

 The coefficient of $y^{m+n-1}$ in $\frac{S_m\cdot S_n}{m\cdot n}$ is $\frac{2}{m}$.

 Therefore, $\boxed{m=2}$.

 Now choose $x=y=\frac1,z=(-2)$. (sic)

 Since $\frac{S_{n+2}}{2+n}=\frac{S_2}{2}\frac{S_n}{n}$ holds for all real $x,y,z$ such that $x+y+z=0$.

 We have $\frac{2^{n+2}-2}{n+2} = 3\cdot\frac{2^n-2}{n}$. Therefore,

 \begin{equation*}\label{eq:l2}\frac{2^{n+1}-1}{n+2} =3\cdot\frac{2^{n-1}-1}{n}\ldots\tag{**}\end{equation*}

 Clearly $(\ref{eq:l2})$ holds for $n\in\{5,3\}$.And one can say that for $n\ge 6$,$\text{RHS of (\ref{eq:l2})}<\text{LHS of (\ref{eq:l2})}$.

 
So our answer is $(m,n)=(5,2),(2,5),(3,2),(2,3)$.

 -ftheftics (edited by integralarefun)

 


	\item (\textit{USAMO 1982}) If a point $A_1$ is in the interior of an equilateral triangle $ABC$ and point $A_2$ is in the interior of $\triangle{A_1BC}$, prove that


 $I.Q. (A_1BC) > I.Q.(A_2BC)$,


 where the \textit{isoperimetric quotient} of a figure $F$ is defined by


 $I.Q.(F) = \frac{\text{Area (F)}}{\text{[Perimeter (F)]}^2}$


 



\textbf{Solution}

First, an arbitrary triangle $ABC$ has isoperimetric quotient (using the notation $[ABC]$ for area and $s = \frac{a + b + c}{2}$):
 \[\frac{[ABC]}{4s^2} = \frac{[ABC]^3}{4s^2 [ABC]^2} = \frac{r^3 s^3}{4s^2 \cdot s(s-a)(s-b)(s-c)} = \frac{r^3}{4(s-a)(s-b)(s-c)}\]
 \[= \frac{1}{4} \cdot \frac{r}{s-a} \cdot \frac{r}{s-b} \cdot \frac{r}{s-c} = \frac{1}{4} \tan A/2 \tan B/2 \tan C/2.\]

 Lemma. $\tan x \tan (A - x)$ is increasing on $0 < x < \frac{A}{2}$, where $0 < A < 90^\circ$.

 Proof. 
 \[\tan x \tan (A - x) = \tan x \cdot \frac{\tan A - \tan x}{1 + \tan A \tan x} = 1 - \frac{1}{\cos^2 x (1 + \tan A \tan x)}\]
 \[= 1 - \frac{2}{1 + \cos 2x + \tan A \sin 2x} = 1 - \frac{2}{1 + \sec A \cos (A - 2x)}\] is increasing on the desired interval, because $\cos (A - 2x)$ is increasing on $0 < x < \frac{A}{2}.$

 Let $x_1, y_1, z_1$ and $x_2, y_2, z_2$ be half of the angles of triangles $A_1 BC$ and $A_2 BC$ in that order, respectively. Then it is immediate that $30^\circ > y_1 > y_2$, $30^\circ > z_1 > z_2$, and $x_1 + y_1 + z_1 = x_2 + y_2 + z_2 = 90^\circ$. Hence, by Lemma it follows that
 \[\tan x_1 \tan y_1 \tan z_1 = \tan (90^\circ - y_1 - z_1) \tan y_1 \tan z_1 > \tan (90^\circ - y_1 - z_2) \tan y_1 \tan z_2\]
 \[> \tan (90^\circ - y_2 - z_2) \tan y_2 \tan z_2 = \tan x_2 \tan y_2 \tan z_2.\]Multiplying this inequality by $\frac{1}{4}$ gives that $I.Q[A_1 BC] > I.Q[A_2 BC]$, as desired.

 


	\item (\textit{USAMO 1982}) Prove that there exists a positive integer $k$ such that $k\cdot2^n+1$ is composite for every positive integer $n$.


 



\textbf{Solution}

Indeed, $\boxed{k=2935363331541925531}$ has the requisite property.

 To see why, consider the primes $3,\ 5,\ 17,\ 257,\ 65537,\ 6700417,\ 641$, and observe that 
 \begin{align*}k&\equiv 1\pmod{3}\\ k&\equiv 1\pmod{5}\\ k&\equiv 1\pmod{17}\\ k&\equiv 1\pmod{257}\\ k&\equiv 1\pmod{65537}\\ k&\equiv 1\pmod{6700417}\\ k&\equiv -1\pmod{641}\end{align*}


 Moreover, 
 \begin{align*}\text{ord}_{3}\left(2\right)&=2\\ \text{ord}_{5}\left(2\right)&=4\\ \text{ord}_{17}\left(2\right)&=8\\ \text{ord}_{257}\left(2\right)&=16\\ \text{ord}_{65537}\left(2\right)&=32\\ \text{ord}_{6700417}\left(2\right)&=64\\ \text{ord}_{641}\left(2\right)&=64\end{align*}


 We conclude that 
 \begin{align*}n\equiv 1\pmod{2}&\implies k\cdot 2^{n}-1\equiv k\cdot\left(-1\right)+1\equiv 0\pmod{3}\\ n\equiv 2\pmod{4}&\implies k\cdot 2^{n}-1\equiv k\cdot\left(-1\right)+1\equiv 0\pmod{5}\\ n\equiv 4\pmod{8}&\implies k\cdot 2^{n}-1\equiv k\cdot\left(-1\right)+1\equiv 0\pmod{17}\\ n\equiv 8\pmod{16}&\implies k\cdot 2^{n}-1\equiv k\cdot\left(-1\right)+1\equiv 0\pmod{257}\\ n\equiv 16\pmod{32}&\implies k\cdot 2^{n}-1\equiv k\cdot\left(-1\right)+1\equiv 0\pmod{65537}\\ n\equiv 32\pmod{64}&\implies k\cdot 2^{n}-1\equiv k\cdot\left(-1\right)+1\equiv 0\pmod{6700417}\\ n\equiv 0\pmod{64}&\implies k\cdot 2^{n}-1\equiv k\cdot1+1\equiv 0\pmod{641}\end{align*}


 And $k>>3,5,17,257,65537,6700417,641$ so the relevant values will, in fact, always be composite. $\blacksquare$

 


	\item (\textit{USAMO 1982}) $A,B$, and $C$ are three interior points of a sphere $S$ such that $AB$ and $AC$ are perpendicular to the diameter of $S$ through $A$, and so that two spheres can be constructed through $A$, $B$, and $C$ which are both tangent to $S$. Prove that the sum of their radii is equal to the radius of $S$.


 



\textbf{Solution}

Let the two tangent spheres be $S_1$ and $S_2$, and let $O, O_1, O_2$ and $R, R_1, R_2$ be the origins and radii of $S, S_1, S_2$ respectively. Then $AO$ stands normal to the plane $P$ through $\Delta ABC$. Because both spheres go through $A$, $B$, and $C$, the line $O_1 O_2$ also stands normal to $P$, meaning $AO$ and $O_1 O_2$ are both coplanar and parallel. Therefore the problem can be flattened to the plane $P'$ through $A$, $O$, $O_1$ and $O_2$. 

 
Let $X, Y, Z, M, N$ be points on $P'$ such that $X = S \cap O O_1,  \quad Y = S_1 \cap S_2 \neq A,  \quad   Z = S \cap O O_2,   \quad  M = O_1 Y \cap AO,    \quad N = O_2 Y \cap AO$ 

 Let $J$ be the three circles radical center, meaning $JX$ and $JZ$ are tangent segments to $S$ and $J \in AY$.

 
Because  $Y \in P \iff \angle NAY = \angle YAM = 90 ^{\circ},$ we have that $\overline{YN}$ and $\overline{YM}$ are diameters. 

 This means that $\angle OZN = \angle ZNY = \angle JZY$ and $\angle MXO = \angle YMX= \angle YXJ$. 

 And because $\angle ZNO = 180^{\circ} - \angle AYZ = \angle ZYJ$ and $\angle OMX = 180^{\circ} -\angle XYA = \angle JYX,$ we have that $\Delta JZY \sim \Delta OZN$ and $\Delta JYX \sim \Delta OMX$. 

 We then conclude that $\overline{OM} = \overline{OX} \enspace \frac {\overline{JY} }{ \overline{JX}}  = \overline{OZ} \enspace \frac {\overline{JY} }{ \overline{JZ}}  = \overline{ON}$

 
Let $a \top b$ denote line $a$ bisecting line segment $b$.

 Since $O_1 O_2 || MN$ it follows that $OY \, \top \, \overline{MN} \Rightarrow OY \, \top \,  \overline{O_1 O_2}$. Similarly we have that $O_1 O_2 \, \top \, \overline{YM} \Rightarrow O_1 O_2 \, \top \,  \overline{OY}$.

 And so $O_1 O O_2 Y$ is a parallelogram because $\overline{OY}$ and $\overline{O_1 O_2}$ bisect each other, meaning $R = \overline{OZ} = \overline{O O_2} + \overline{O_2 Z} = \overline{O_1 Y}+ R_2 = R_1 + R_2$

 
$Quod \enspace Erat \enspace Demonstrandum$

 


\end{enumerate}